% Self-contained appendix input. Requires amsmath, amssymb only.
% No document class, bibliography commands, or project-specific macros.
\section{Conditional EI analytics and the population geometry}
\label{app:final_conditional_analytics}

This appendix completes the algebra of the conditional construction used in
this paper. The construction is an asymptotic cavity closure. Exact
finite-dimensional identities, algebraic implications of that closure, and
additional assumptions about the nonlinear network are kept separate. The
finite-network tests determine where the resulting parameter-only prediction
is quantitatively useful.

\subsection{Model, averaging operations, and the observable}

Let $i\in a$ identify a neuron in population $a$, with $N_a=f_aN$ and
$\sum_a f_a=1$. Set $F=\operatorname{diag}(f_a)$. In common leak-time units,
\begin{equation}
 \dot x_i=-x_i+\sum_jJ_{ij}\phi(x_j)+\sqrt N I_a+h_a^{\rm ext},
 \qquad \phi(x)=2\Phi(x),\qquad r_i=\phi(x_i).
 \label{eq:an_model}
\end{equation}
The dense Bernoulli ensemble has $J_{ij}=t_{ab}A_{ij}/\sqrt N$,
$t_{aE}>0$, $t_{aI}<0$, and independent $A_{ij}$ with success probability
$p_{ab}$. Its moment coefficients are
$J^0_{ab}=p_{ab}t_{ab}$ and $G_{ab}=p_{ab}(1-p_{ab})t_{ab}^2$.
Thus $\mathbb E J_{ij}=J^0_{ab}/\sqrt N$ and
$\operatorname{Var}J_{ij}=G_{ab}/N$. Self-connections follow the same rule.
Degree-constrained rows are treated below.

Time averaging within a stationary state is denoted by angle brackets;
$\mathbb E_z$ averages fixed neuronal operating points, and an ensemble average
over connectivity realizations is a further operation. Define
\begin{align}
 m_i&=\langle r_i\rangle,&
 C_{ij}(\tau)&=\langle[r_i(t+\tau)-m_i][r_j(t)-m_j]\rangle,\\
 \nu_N&=\frac{\operatorname{tr}C(0)}N,&
 \mathcal E_N&=\frac{\operatorname{tr}[C(0)^2]}N,&
 d_N&=\frac{D_{{\rm PR},N}}N=\frac{\nu_N^2}{\mathcal E_N}.
 \label{eq:an_observables}
\end{align}
At zero lag $C$ is symmetric, so
$\operatorname{tr}C^2=\sum_iC_{ii}^2+\sum_{i\ne j}C_{ij}^2$.
Participation ratio therefore depends on the second trace moment, rather
than on the rank or on the Lyapunov spectrum. At a true fixed point the
time-centered covariance vanishes and this ratio is undefined.

\subsection{The conditional one-site law}

Cancellation of the $O(\sqrt N)$ mean recurrent and external currents gives
\begin{equation}
 J^0F\overline m+I=0.
 \label{eq:an_balance}
\end{equation}
Assume this equation selects admissible mean rates. The conditional cavity law
is
\begin{equation}
 x_a(t)\mid z=\mu_a(z)+u_a(t),\qquad
 \mu_a(z)=h_a+\sqrt{\Delta_a}z,\qquad
 z\sim\mathcal N(0,1),\quad
 u_a\sim\operatorname{GP}(0,K_a),\quad q_a=K_a(0).
 \label{eq:an_site}
\end{equation}
The quenched $z$ remains fixed along a neuron's trajectory. The temporal
Gaussian process is independent of $z$ in this closure. Its entire kernel
$K_a$ is part of the law; its variance alone does not specify it. The
nonlinear rate process is generally non-Gaussian.

Conditional temporal means and averaged slopes follow by integrating two
Gaussian variables:
\begin{align}
 m_a(z)&=2\Phi\!\left(\frac{\mu_a(z)}{\sqrt{1+q_a}}\right),\\
 \alpha_a(z)&=\mathbb E_u\phi'(\mu_a(z)+u)
 =\sqrt{\frac2{\pi(1+q_a)}}
   \exp\!\left[-\frac{\mu_a(z)^2}{2(1+q_a)}\right],\\
 \overline m_a&=2\Phi\!\left(\frac{h_a}{\sqrt{1+\Delta_a+q_a}}\right),
 \qquad
 h_a=\sqrt{1+\Delta_a+q_a}\,
       \Phi^{-1}(\overline m_a/2).
 \label{eq:an_site_closed}
\end{align}
For example, $m(\mu)=2\mathbb P(Y\leq\mu+\sqrt q Z)$ for independent standard
normal $Y,Z$; their difference has variance $1+q$, giving the first formula.
Completing the square in the corresponding Gaussian density gives $\alpha$.

For fixed $\mu,q$, define the centered Gaussian rate covariance
$\mathcal C(\mu,q;k)$ when the two inputs have covariance $k$. Price's theorem
and independence at $k=0$ give
\begin{align}
 \partial_k\mathcal C(\mu,q;k)
 &=\frac{2\exp[-\mu^2/(1+q+k)]}
         {\pi\sqrt{(1+q)^2-k^2}},\\
 \mathcal C(\mu,q;k)
 &=\frac2\pi\int_0^{\arcsin[k/(1+q)]}
 \exp\!\left[-\frac{\mu^2}{(1+q)(1+\sin\theta)}\right]d\theta.
 \label{eq:an_price}
\end{align}
Consequently $c_a(z,\tau)=\mathcal C(\mu_a(z),q_a;K_a(\tau))$ and
$v_a(z)=c_a(z,0)$. The static and dynamic self-consistency equations are
\begin{align}
 \Delta_a&=\sum_bG_{ab}f_b\mathbb E_z[m_b(z)^2],\\
 C_a^\eta(\tau)&=\sum_bG_{ab}f_b\mathbb E_z[c_b(z,\tau)],&
 K_a''(\tau)&=K_a(\tau)-C_a^\eta(\tau),
 \label{eq:an_selfconsistency}
\end{align}
with $K_a'(0)=0$, $K_a(0)=q_a$, and $K_a(\infty)=0$ on the stationary mixing
branch. The first equation contains the second moment of the conditional
mean, $\mathbb E m_b^2=\overline m_b^2+\operatorname{Var}_z m_b$. Heterogeneity
survives the dense limit because this quenched input variance remains finite.

Use the transform $c(\omega)=\int e^{-i\omega\tau}c(\tau)d\tau$ and inverse
measure $d\omega/(2\pi)$. The leak filter is
\begin{equation}
 H(\omega)=\frac1{1+i\omega},\qquad
 K_a(\omega)=|H(\omega)|^2 C_a^\eta(\omega),\qquad
 s_a(z,\omega)=\alpha_a(z)H(\omega).
 \label{eq:an_filter}
\end{equation}
Here $s$ is the time-averaged local response to additive current before
recurrent propagation. It is not the pathwise tangent propagator of a chaotic
network, and replacing an entire time-dependent tangent matrix by averaged
slopes is not an exact stationary-response identity.

The Hermite expansion makes the nonlinear part explicit. Write
\[
 \phi(\mu+\sqrt q Z)-m
 =\sum_{n\geq1}b_n(\mu)\frac{\operatorname{He}_n(Z)}{\sqrt{n!}}.
\]
Gaussian integration by parts gives
\begin{equation}
 b_n(\mu)=\frac{q^{n/2}}{\sqrt{n!}}
          \mathbb E\phi^{(n)}(\mu+\sqrt q Z),\qquad
 c(\mu,\tau)=\sum_{n\geq1}b_n(\mu)^2[K(\tau)/q]^n.
 \label{eq:an_hermite}
\end{equation}
In particular $b_1=\sqrt q\alpha$, and
$c(\mu,\omega)-\alpha(\mu)^2K(\omega)$ is a nonnegative spectrum. The
same label controls variance, spectrum, and response. Their joint moments
must be retained.

\subsection{From the two-site closure to a small population resolvent}

The additional cavity assumptions concern correlations between distinct
neurons and their reservoir: weak pair covariances of order $N^{-1/2}$,
conditional decoupling of the required reservoir response contractions, exact
retention of repeated-site conditional moments, and bounded source-weighted
response norms at fixed parameters. These are stronger than the one-site
Gaussian law. They are not implied by it.

At each pair of frequencies define
\begin{equation}
 \begin{aligned}
 A_a&=\mathbb E_z\alpha_a(z)^2,\qquad
 &u_a&=A_aH(\omega_1)H(\omega_2),\qquad
 &U&=\operatorname{diag}(u_a),\\
 Q_a&=\mathbb E_z[c_a(z,\omega_1)c_a(z,\omega_2)],
 &Q_{\rm d}&=\operatorname{diag}(Q_a).
 \end{aligned}
 \label{eq:an_AUQ}
\end{equation}
Both factors in $Q_a$ and $u_a$ belong to the same neuron. For distinct
$i\in a,j\in b$, let
$\Psi_{ab}=\lim_{N\to\infty}N\mathbb E[C_{ij}(\omega_1)C_{ij}(\omega_2)]$
under the closure. The shared-input and directed-interaction contractions
satisfy
\begin{align}
 \Gamma_\eta&=GFQ_{\rm d}G^T+GF\Psi FG^T,\\
 \Gamma_F&=G+GFU^*G+GFU^*\Gamma_FU^*FG,\\
 \Psi&=U\Gamma_\eta U^*+U\Gamma_F^*Q_{\rm d}
                   +Q_{\rm d}\Gamma_F^TU^*.
 \label{eq:an_cavity}
\end{align}
A star denotes complex conjugation, and a dagger conjugate transpose. The
last two terms retain directed propagation through the reservoir; omitting
them is not the two-site closure used here.

For the regular branch $\Gamma_F=(I-GFU^*)^{-1}G$. Set
\begin{equation}
 L=UGF,\qquad R=(I-L)^{-1},\qquad
 \mathsf D=\operatorname{diag}(Q_a/f_a),\qquad X=\Psi+\mathsf D.
 \label{eq:an_R}
\end{equation}
Substitution in the last cavity equation gives
\begin{align}
 \Psi&=L(\Psi+\mathsf D)L^\dagger
                  +RL\mathsf D+\mathsf DL^\dagger R^\dagger,\\
 X&=LXL^\dagger+R\mathsf D+\mathsf DR^\dagger-\mathsf D,
 \label{eq:an_Lyapunov}
\end{align}
where $RL=R-I$. To verify the solution, use $LR=R-I$ on both sides:
\begin{align}
 L(R\mathsf DR^\dagger)L^\dagger
  &=(R-I)\mathsf D(R^\dagger-I)\\
  &=R\mathsf DR^\dagger-R\mathsf D-\mathsf DR^\dagger+\mathsf D.
\end{align}
Adding the remaining terms gives $X=R\mathsf DR^\dagger$. If $\rho(L)<1$,
the discrete Lyapunov equation has a unique solution, hence
\begin{equation}
 \boxed{X=R\mathsf DR^\dagger,\qquad \Psi=X-\mathsf D.}
 \label{eq:an_final_pair}
\end{equation}
This implication is exact algebra conditional on Eq.~\eqref{eq:an_cavity}.
For nonnegative $G$, a sufficient condition at all real frequencies is
$\rho(B)<1$, where $B=\operatorname{diag}(A_a)GF$.

This convention fixes the lag orientation independently of the derivation.
An isolated $j\in b\to i\in a$ edge gives
$C_{ij}(\omega)=J_{ij}s_i(\omega)c_j(\omega)$ to first order, so its
contribution is $G_{ab}u_aQ_b$, rather than $G_{ab}u_a^*Q_b$.
The inverse transform of $\Psi$ uses
$e^{i\omega_1\tau_1+i\omega_2\tau_2}$. It obeys
$\psi_{ba}(\tau_1,\tau_2)=\psi_{ab}(-\tau_1,-\tau_2)$ and
$\psi_{ab}(\tau_1,\tau_2)=\psi_{ab}(\tau_2,\tau_1)$.
Conjugating the response resolvent reverses the lag map, while leaving the
zero-lag energy unchanged.

\subsection{Population covariance energies and participation ratios}

Define
\begin{equation}
 S_{ab}=\iint X_{ab}\frac{d\omega_1d\omega_2}{(2\pi)^2},\qquad
 \psi_{ab}=\iint\Psi_{ab}\frac{d\omega_1d\omega_2}{(2\pi)^2},\qquad
 V_{2a}=\mathbb E_zv_a(z)^2,\quad \overline v_a=\mathbb E_zv_a(z).
\end{equation}
Since $\iint Q_a=V_{2a}$, Eq.~\eqref{eq:an_final_pair} gives
$S_{ab}=\psi_{ab}+\delta_{ab}V_{2a}/f_a$. The physical block energies,
normalized by the total neuron count, are
\begin{equation}
 \mathcal E_{ab}
 =\lim_{N\to\infty}\frac1N
           \sum_{i\in a,j\in b}C_{ij}(0)^2
 =f_af_bS_{ab}
 =f_af_b\psi_{ab}+\delta_{ab}f_aV_{2a}.
 \label{eq:an_blocks}
\end{equation}
The off-diagonal pairs each contribute $O(N^{-1})$ but number $O(N^2)$;
they therefore contribute $O(1)$ to $\mathcal E$. Both orientations of the
E--I block are included in the total:
\begin{equation}
 \mathcal E=\mathcal E_{EE}+2\mathcal E_{EI}+\mathcal E_{II}
           =f^TSf,\qquad
 \nu=\sum_af_a\overline v_a,\qquad
 \boxed{d=\frac{D_{\rm PR}}N=\frac{\nu^2}{\mathcal E}.}
 \label{eq:an_dimension}
\end{equation}
Within one population the correct normalization is
\begin{equation}
 \boxed{\frac{D_{{\rm PR},a}}{N_a}
    =\frac{\overline v_a^2}{V_{2a}+f_a\psi_{aa}}
    =\frac{\overline v_a^2}{f_aS_{aa}}.}
 \label{eq:an_restricted}
\end{equation}
For example, $D_{{\rm PR},E}/N$ has an additional factor $f_E$.

With $\mathcal E_{\rm diag}=\sum_af_aV_{2a}$, the useful physical
factorization is
\begin{equation}
 d=\underbrace{\frac{(\sum_af_a\overline v_a)^2}
                       {\sum_af_aV_{2a}}}_{d_{\rm uncorrelated}}
   \underbrace{\frac{\mathcal E_{\rm diag}}{\mathcal E}}_{
                    \text{reduction by pair correlations}}.
 \label{eq:an_factorization}
\end{equation}
The first factor is at most one by Jensen and Cauchy--Schwarz. It is the PR
fraction of independent neurons having the same distributed variances.
The second factor measures the additional effect of pair correlations.
The analogous within-population factors are
$\overline v_a^2/V_{2a}$ and $V_{2a}/(V_{2a}+f_a\psi_{aa})$.

For nonnegative conditional spectra, $X(\omega_1,\omega_2)$ is Hermitian
positive semidefinite. The integrated $S$ is real symmetric and positive
semidefinite. In particular,
\begin{equation}
 0\leq\frac{\mathcal E_{EI}}
              {\sqrt{\mathcal E_{EE}\mathcal E_{II}}}\leq1,
 \qquad
 0\leq\frac{2\mathcal E_{EI}}{\mathcal E}\leq\frac12,
 \label{eq:an_block_bounds}
\end{equation}
where zero-denominator cases are omitted. Nonnegativity of the integrated
entries also follows from the path expansion below. These are second-moment
properties of population covariance geometry. An E--I energy is not a PR of a
rectangular matrix, a communication dimension, or a determination of the
singular vectors of the cross-covariance block.

\subsection{Premature averaging inequalities for a nonuniform population profile}
\label{app:population_averaging_inequalities}

The sign of the averaging error can be established for the full population
problem with a common leak filter, not just the uniform scalar reduction.
Assume $B\geq0$, $\rho(B)<1$, nonnegative integrable conditional spectra,
and finite $V_{2a}$. Writing $t=H(\omega_1)H(\omega_2)$ gives
$R=\sum_{n\geq0}t^nB^n$. For each conditional law define the real number
\begin{equation}
 a_{nm,c}(z)=\int c_c(z,\omega)H(\omega)^nH(\omega)^{*m}
                   \frac{d\omega}{2\pi}.
 \label{eq:an_real_moment}
\end{equation}
Reality follows by changing $\omega\to-\omega$ in its complex conjugate.
Moreover $|a_{nm,c}(z)|\leq v_c(z)$ because $|H|\leq1$.
Substituting the resolvent series into $S$ and using the same label at both
frequencies yields
\begin{equation}
 \boxed{S_{ab}=\sum_{c,n,m\geq0}
   \frac{(B^n)_{ac}(B^m)_{bc}}{f_c}
          \mathbb E_z[a_{nm,c}(z)^2].}
 \label{eq:an_population_positive_series}
\end{equation}
The domination by $V_{2c}$ and the convergent nonnegative matrix geometric
series justify interchanging the sums and integrals. Every term is
nonnegative.

First replace $Q_c$ by the product of population-averaged spectra, while
holding $A$ and each $\overline v_c$ fixed. The source second moment in
Eq.~\eqref{eq:an_population_positive_series} becomes
$[\mathbb E_z a_{nm,c}(z)]^2$. Thus
\begin{equation}
 S_{ab}-S_{ab}^{\rm average\ spectra}
 =\sum_{c,n,m\geq0}\frac{(B^n)_{ac}(B^m)_{bc}}{f_c}
                 \operatorname{Var}_z[a_{nm,c}(z)]\geq0.
 \label{eq:an_Q_inequality}
\end{equation}
Next replace $A_c=\mathbb E\alpha_c^2$ by
$A_c^{\rm average\ gain}=(\mathbb E\alpha_c)^2$, holding the conditional
spectra fixed. Jensen gives $0\leq B^{\rm average\ gain}\leq B$
entrywise. Every matrix power has the same ordering, so every $S_{ab}$
can only decrease. The unchanged numerator then implies that both total and
population-restricted PR fractions can only increase in either approximation.
These statements compare approximations to the same one-site law; they do
not compare distinct self-consistent network interventions. They also do not
claim pointwise monotonicity of complex frequency-resolved matrix entries.

In the reference nonuniform law, the resulting PR fractions are
\begin{center}
\begin{tabular}{c|rrrr}
$f_E$ & Conditional & Average spectra & Average gain & Both\\\hline
$0.6$ & $0.0298935$ & $0.0665594$ & $0.163924$ & $0.359971$\\
$0.8$ & $0.0343262$ & $0.0764508$ & $0.168014$ & $0.368452$
\end{tabular}
\end{center}
These are parameter-only quadratures at fixed aggregate coupling and fixed
conditional laws. All blockwise inequalities pass independently; the
conditional baselines reproduce the previous calculation to floating
precision. No finite-network observation was used to choose these values.

\subsection{Composition and degree controls change different objects}

At fixed target mean rates, holding $W=GF$ and $J^0F$ fixed preserves the
conditional one-site equations. When $f$ changes, rescale each presynaptic
column as $G_{ab}^{\rm new}=G_{ab}^{\rm ref}f_b^{\rm ref}/f_b^{\rm new}$,
and apply the same rule to $J^0$. Then $R$ and $Q$ remain fixed, but
$\mathsf D_{cc}=Q_c/f_c$ and the external block weights change. The
fraction-independent tensor
\begin{equation}
 T_{abc}=\iint R_{ac}R_{bc}^*Q_c
                  \frac{d\omega_1d\omega_2}{(2\pi)^2}
 \quad\hbox{gives}\quad S_{ab}(f)=\sum_c\frac{T_{abc}}{f_c}.
 \label{eq:an_composition}
\end{equation}
Keeping individual edge moments fixed is a different control path requiring
new one-site laws. Identical local laws with uniform edge variance instead
give the null result $\mathcal E=\iint Q/|1-g^2AH_1H_2|^2$, independent of
population fractions.

For the degree intervention, keep $J^0,G,f$ and the target mean rates fixed,
choose an incoming count $K_{iab}$, and uniformly sample a subset of the
$M=N_b$ possible presynaptic neurons. Define
$\chi_{ab}=\operatorname{Var}K_{iab}/[Mp_{ab}(1-p_{ab})]$.
For fixed presynaptic values with mean $\overline m$ and variance $V$, exact
sampling without replacement gives
\begin{equation}
 \operatorname{Var}\!\left(\frac{t}{\sqrt N}\sum_jA_jm_j\right)
 =\frac{t^2}{N}Mp(1-p)
     \left[\frac{M-\chi}{M-1}V+\chi\overline m^2\right].
 \label{eq:an_finite_degree}
\end{equation}
For bounded $\chi=O(1)$, the leading static law is therefore
\begin{equation}
 \Delta_a=\sum_bG_{ab}f_b
  \left[\operatorname{Var}_z m_b(z)+\chi_{ab}\overline m_b^2\right].
 \label{eq:an_degree_static}
\end{equation}
The finite-$M$ correction and the actual generator dispersion can be retained
as diagnostics. The dynamic and pair equations retain their form only under
the dense cavity assumptions, including control of coherent population
projections. Degree constraints change within-row covariances; their
population-direction part is not zero merely because the mean EI matrix is
Hurwitz. Bounded source-weighted response operators and subextensive
coherent-projection covariance energy are additional conditions. Thus degree
variation changes the local laws, $A$, and $Q$ self-consistently, whereas the
fixed-$W$ composition control reweights the same laws.

\subsection{Coherent EI feedback and endogenous labels}

The mean balanced connectivity must remain in the dynamics. Write
$J=J^{\rm c}+P K_0 B_0$, with $P_{ia}=1_{i\in a}$,
$(B_0)_{aj}=1_{j\in a}/N_a$, and $K_0=\sqrt N J^0F$.
For a frozen local-response comparison let
$L_0=\operatorname{diag}(s_i)$ and $M_0=(I-L_0J^{\rm c})^{-1}$.
Where the inverses exist, Woodbury gives the exact identity
\begin{equation}
 (I-L_0J)^{-1}=M_0+
 M_0L_0P(I-K_0B_0M_0L_0P)^{-1}K_0B_0M_0.
 \label{eq:an_coherent}
\end{equation}
It retains the full coherent EI loop. This finite-dimensional comparison is
not an identity for a trajectory-dependent nonlinear tangent propagator.
Removing the loop before solving for operating points is especially
unjustified: it changes the one-site law and can destroy stability.

Even conditioning on a neuron's quenched offset changes its incoming row.
Hold the presynaptic mean-rate vector $m_j$ fixed in a cavity reservoir,
independent of the added neuron's incoming Gaussian row. For independent
Gaussian $J_{ij}^{\rm c}$ of variances $G_{ab}/N$, set
$\eta_i=\sum_jJ_{ij}^{\rm c}m_j$, $v_j=G_{ab}m_j/N$, and
$\sigma_i^2=\sum_jG_{ab}m_j^2/N$. Exact Gaussian regression yields
\begin{equation}
 \mathbb E[J_{ij}^{\rm c}\mid\eta_i]=\eta_i v_j/\sigma_i^2,
 \qquad
 \operatorname{Cov}(J_{ij}^{\rm c},J_{ik}^{\rm c}\mid\eta_i)
   =\frac{G_{ab}}N\delta_{jk}-\frac{v_jv_k}{\sigma_i^2}.
 \label{eq:an_row_regression}
\end{equation}
The conditional mean and removed direction produce finite-rank structure.
The fixed-reservoir premise matters: the actual recurrent $m_j$ depends on
the full connectivity, so the Gaussian prior alone does not make this an
exact regression formula for the full nonlinear network. A Bernoulli
counterpart additionally requires a conditional central-limit approximation.
Finite rank by itself does not bound a covariance correction near an
amplified mode. In a linear residual comparison with fixed source spectra,
uniformly bounded source-weighted resolvents and an integrable envelope make
a rank-$r$ change contribute $O(r)$ to the trace norm of the integrated
covariance difference. If covariance operator norms are also bounded,
\begin{equation}
 |\operatorname{tr}(C-\widetilde C)|=O(r),\qquad
 |\operatorname{tr}(C^2-\widetilde C^2)|
 \leq\|C-\widetilde C\|_1\|C+\widetilde C\|_{\rm op}=O(r).
\end{equation}
Its normalized trace moments are then subleading for fixed $r$. Re-solving
the nonlinear network can change the source in full rank; the comparison
does not prove that change small. This is why the coherent and conditional
cavity assumptions must be tested separately in finite networks.

There is also an exact finite-network balance correction. Averaging the
stationary input equation over each population gives
\begin{equation}
 \overline m_N-\overline m_*=
 N^{-1/2}(J^0F)^{-1}
   (\overline\mu_N-h^{\rm ext}-\overline\eta_N),
 \qquad \overline m_*=-(J^0F)^{-1}I.
 \label{eq:an_finite_balance}
\end{equation}
The centered recurrent static current $\overline\eta_N$ must be retained.
Dropping it is not an exact finite-gain balance constraint.

\subsection{Critical scope and checks}

A critical reduction additionally requires an isolated response eigenmode,
a smooth stationary branch, stable coherent feedback, and nonzero projected
quadratic curvature. The local covariance expansion is
$c_a(z,\tau)=\alpha_a(z)^2K_a(\tau)+\beta_a(z)^2K_a(\tau)^2/2+\cdots$,
where $\beta_a=\mathbb E_u\phi''(\mu_a+u)$; averaging retains
$\mathbb E_z\beta_a^2$, rather than $(\mathbb E_z\beta_a)^2$.
For the generic quadratic scalar reduction, with
$\delta=1-g^2\mathbb E\alpha^2$ the scalar response gap,
$K''=\delta K-g_c^2\mathbb E\beta^2 K^2/2+\cdots$ and the zero-energy
turning point gives $q=3\delta/(g_c^2\mathbb E\beta^2)+\cdots$.
The familiar pair resonance then produces the conditional $d\propto
\delta^{3/2}$ law. Degenerate modes or vanishing quadratic curvature require
a different reduction. A biased homogeneous operating point can supply
curvature, so the exponent alone does not identify a heterogeneity mechanism.
No finite-network critical exponent is established by this algebra.

Independent checks accompanying this appendix compare Price integrals with
two-dimensional Gaussian quadrature, the matrix closure with a separate
Lyapunov solver, the population path series with direct quadrature, the
uniform-fraction null limit, uncoupled diagonal energy, and an isolated
feedforward edge's phase. Maximum absolute algebra errors are below
$4.1\times10^{-14}$ and the path-series relative error is below
$1.8\times10^{-15}$ in the supplied cases. These checks verify integrals,
normalizations, and algebra. They do not validate the nonlinear cavity
assumptions or replace network sampling.
